\documentclass[12pt, a4paper]{article}

% ====== Packages ======
\usepackage[utf8]{inputenc}
\usepackage[T1]{fontenc}
\usepackage{amsmath, amssymb, amsfonts}
\usepackage{geometry}
\geometry{top=2.5cm, bottom=2.5cm, left=2.5cm, right=2.5cm}
\usepackage{graphicx}
\usepackage{booktabs}
\usepackage{tabularx}
\usepackage{xcolor}
\usepackage{hyperref}
\usepackage{listings}
\usepackage{enumitem}

% ====== Hyperref Setup ======
\hypersetup{
    colorlinks=true,
    linkcolor=blue,
    filecolor=magenta,      
    urlcolor=cyan,
    pdftitle={In-Depth Design Analysis Report: MMHCL Framework Upgrade},
}

% ====== Code Listing Setup ======
\definecolor{codegreen}{rgb}{0,0.6,0}
\definecolor{codegray}{rgb}{0.5,0.5,0.5}
\definecolor{codepurple}{rgb}{0.58,0,0.82}
\definecolor{backcolour}{rgb}{0.95,0.95,0.92}

\lstdefinestyle{mystyle}{
    backgroundcolor=\color{backcolour},   
    commentstyle=\color{codegreen},
    keywordstyle=\color{magenta},
    numberstyle=\tiny\color{codegray},
    stringstyle=\color{codepurple},
    basicstyle=\ttfamily\footnotesize,
    breakatwhitespace=false,         
    breaklines=true,                 
    captionpos=b,                    
    keepspaces=true,                 
    numbers=left,                    
    numbersep=5pt,                  
    showspaces=false,                
    showstringspaces=false,
    showtabs=false,                  
    tabsize=4
}
\lstset{style=mystyle}

% ====== Document Meta ======
\title{\textbf{In-Depth Design Analysis Report:\\Upgrading the MMHCL Multi-Modal Recommendation Framework via Spectral Domain Representation Calibration (DAMPS)}}
\author{\textit{Revision 10 — Phase 1 (Quick Win) Execution Plan}}
\date{}

\begin{document}

\maketitle

\section*{Executive Summary}
In the booming landscape of e-commerce platforms, Recommender Systems are increasingly facing the immense challenge of data sparsity. The Multi-Modal Hypergraph Contrastive Learning (MMHCL) framework pioneered a novel direction by modeling shared preferences through hypergraphs to enrich feature representations. 

Despite its superior performance, an in-depth architectural analysis reveals that the item-item (i2i) hypergraph construction mechanism of MMHCL harbors a core vulnerability: the direct use of raw multi-modal features to compute similarities causes the network to be heavily dominated by noise. Empirical measurements indicate that up to 44\% of visual features and 23\% of textual features are dominated by junk noise.

The proposed solution directly integrates the Adaptive Phase and Spectral Calibration (DAMPS) framework into the MMHCL architecture. However, the latest empirical analysis on the Amazon Clothing dataset demonstrates that the model is experiencing an asymmetric coverage loss phenomenon, where Recall@20 drops by 10.7\% while NDCG@20 only decreases by 2.9\% compared to the original baseline. To thoroughly address this issue, the current blueprint restricts its focus to \textbf{Phase 1 — Quick Win ($<$ 1 day)} with 3 core actions: \textbf{Audit eval protocol}, \textbf{Static $\tau$ sweep}, and \textbf{Amplitude Variance Ratio Filter (AVRF) ablation} for sparse data. The combined execution of Phase 1 is expected to yield a recovery of +5\%\textasciitilde8\% in Recall@20, assuming the evaluation protocol perfectly aligns. More advanced phases (Phase 2, Phase 3) are reserved and will be deployed subsequently.

\section{Research Context and Mathematical Flaws of MMHCL}
To understand why this upgrade is imperative, we must analyze the current mathematical structure of MMHCL during the hypergraph initialization process.

\subsection{Current i2i Hypergraph Construction Mechanism}
The objective of MMHCL is to link items with semantic similarities through a graph structure. Initially, the system computes the Cosine similarity matrix ($I_{ij}^{m}$) directly from the raw feature vectors $e_i^m$ and $e_j^m$ of item $i$ and $j$ in modality $m$ (image/text):
\begin{equation}
I_{ij}^{m} = \frac{(e_i^m)^T e_j^m}{||e_i^m|| \cdot ||e_j^m||}
\end{equation}
After obtaining the dense similarity matrix, MMHCL applies the K-Nearest Neighbors (K-NN) algorithm to filter out noise, retaining only the hyperedges with the highest weights to construct a sparse matrix $\hat{I}_{ij}^{m}$:
\begin{equation}
\hat{I}_{ij}^{m} = \begin{cases} 1, & I_{ij}^{m} \in \text{top-K}(I_i^m) \\ 0, & \text{otherwise} \end{cases}
\end{equation}
Finally, the uni-modal graphs are fused to form the complete hypergraph.

\subsection{Core Architectural Flaw Analysis}
From a signal processing perspective, calculating the Cosine similarity directly on the raw vectors $e_i^m$ without passing them through a noise filter poses a significant risk. E-commerce data inherently contains substantial localized noise (e.g., cluttered image backgrounds, redundant text descriptions).

Consequently, instead of reflecting ``semantic distance'', the dot product in the Cosine formula captures the similarity of ``noise vectors''. When the K-NN algorithm selects the nearest neighbors, it inadvertently picks items with identical noise patterns, creating \textbf{phantom edges}. This pollutes the downstream contrastive learning embedding space, causing severe performance degradation, especially as the coefficient $K$ increases.

\subsection{DAMPS Integration: Zero-Risk Solution with Soft Residual-Routing}
The intervention strategy involves inserting the DAMPS filter between the feature extraction and Cosine computation steps. Instead of the conventional ``Dual-Routing'' approach—which allows 44\% of raw noise to leak through the GCN network—the design adopts a \textbf{Soft Residual-Routing} model.

The input features for the HGCN are established in a residual format:
\begin{equation}
h_{input} = h_{raw} + \alpha \cdot \text{LayerNorm}(h_{cal})
\end{equation}
where $\alpha$ is a learnable parameter (initialized at 0.1). This mechanism allows DAMPS to partially clean the node features, preventing raw noise from propagating through the GCN while preserving personalized information diversity.

\section{Mathematical Foundation and Analysis of DAMPS Modules}
The superiority of the proposed method lies in decoupling multi-modal representations into two independent components: the Amplitude Spectrum, indicating energy magnitude, and the Phase Spectrum, encoding relative spatial shifts.

\subsection{Multi-Modal Spectral Decomposition}
The initial raw uni-modal features $x_i^m$ are mapped by MMHCL through MLP projection layers:
\begin{equation}
h_i^m = W_m x_i^m + b_m
\end{equation}
The algorithm applies the Fast Fourier Transform (FFT) on the representations $h_i^m$, transforming them into complex numbers in the spectral domain $z_i^m \in \mathbb{C}^F$, with the number of frequency bins $F = \frac{d}{2} + 1$. At each frequency bin $f$, the signal is mathematically represented in polar form:
\begin{equation}
z_{i,f}^m = |z_{i,f}^m|e^{j\phi_{i,f}^m}
\end{equation}
where $|z_{i,f}^m|$ is the amplitude and $\phi_{i,f}^m$ is the phase.

\subsection{Metadata-Aware Phase Calibrator (Metadata-Aware APC)}
The phase distribution shift between images and texts artificially narrows the similarity computations $I_{ij}^{m}$. Utilizing the dynamic $k$-means clustering algorithm creates a severe CPU bottleneck due to continuous CPU-GPU synchronization.

To overcome this, the system is optimized into a \textbf{Metadata-Aware APC}, relying on static product metadata categories for clustering. Subsequently, von Mises distribution Maximum Likelihood Estimation (MLE) is applied to compute an independent central phase rotation angle $\hat{\theta}_c$ for each cluster $c$:
\begin{equation}
\hat{\theta}_c = \text{atan2}\left(\sum_{i \in C_c} \sin R_i, \sum_{i \in C_c} \cos R_i\right)
\end{equation}
Then, a trainable residual rotation parameter $\psi$ is applied:
\begin{equation}
\tilde{z}_i^{image} = z_i^{image}e^{-j(\frac{\hat{\theta}_c}{2} + \psi)}, \quad \tilde{z}_i^{text} = z_i^{text}e^{+j(\frac{\hat{\theta}_c}{2} + \psi)}
\end{equation}

\subsection{Amplitude Variance Ratio Filter (AVRF) and Phase 1 Ablation Plan}
To address the 44\% noise volume within the amplitude spectrum, AVRF utilizes the Median Absolute Deviation (MAD). The filter is designed in a compact form inspired by the Wiener filter:
\begin{equation}
V^m(\tilde{A}_{I,f}^m) = \frac{\sigma_{inter}^2}{\sigma_{inter}^2 + \sigma_{intra}^2 + \epsilon}
\end{equation}
\textbf{Phase 1 Evaluation (AVRF Ablation):} Although AVRF possesses strong noise filtering capabilities, the latest review panel analysis indicates that on a sparse dataset like Amazon Clothing, the AVRF spectral filter tends to \textit{over-attenuate} useful signals. This occurs because the raw features (e.g., 4096-D images, text from pretrained encoders) have already been partially denoised beforehand. Therefore, within the execution framework of \textbf{Phase 1 — Quick Win}, AVRF will be temporarily disabled (ablated) to recover dataset coverage and prevent the loss of sensitive signals in the long-tail region.

\subsection{Inter-Modal Coherence Filter (Residual IMCF Form)}
Empirical measurements indicate that the Amplitude Spectral Coherence (ASC) coefficient on e-commerce data generally fluctuates at a low level (mode $0.2-0.4$). Directly multiplying this value would cause severe multiplicative attenuation. To resolve this, IMCF is restructured into a residual form, modulating only the values around the baseline $\overline{ASC}$:
\begin{equation}
z_i^{m*} = \tilde{z}_i^{m(vrf)} + \lambda_{coh}(ASC_i - \overline{ASC}) \odot \tilde{z}_i^m
\end{equation}
The system then utilizes Inverse FFT (IFFT) to project the cleaned data back into the original spatial domain.

\section{System Integration Design and Training Dynamics}
The transition from theory to practical deployment on an RTX 5090 requires rigorous architectural adjustments to avoid convergence cliffs and memory explosions.

\subsection{Pattern B' (Scheduled Rebuild) and Dynamics}
Invoking the DAMPS filter and recomputing the K-NN hypergraph at every forward pass creates instability. Therefore, the system adopts \textbf{Pattern B'}, with a rebuild frequency of every $R$ epochs (e.g., $R=5$).

\noindent \textbf{Fixing the Density Explosion Bug:}\\
The system is mandated to use a Momentum Encoder directly on the embedding representation space itself to maintain a stable number of connections, preventing VRAM overflow (OOM).

\noindent \textbf{Static InfoNCE Temperature (Static $\tau$ Sweep) for Phase 1:}\\
Previously, the Momentum Encoder smoothed representations and required $\tau$ to be a learnable parameter. However, empirical analysis revealed that the gradient of $\tau$ vanished, leaving its value permanently stuck at 0.0909 across all 10 seeds. This directly triggered an \textit{embedding collapse} and severely degraded coverage (evidenced by the 10.7\% plunge in Recall@20). For immediate remediation in the Phase 1 campaign, the system pivots to using a \textbf{Static $\tau$ sweep} with candidate sets $\{0.2, 0.3, 0.5\}$, prioritizing a fixed anchor at $\tau = 0.3$.

\noindent \textbf{Training Diagnostics:}\\
To ensure transparency, the system implements two diagnostic metrics inside the training log:

\begin{lstlisting}[language=Python, caption=Proposed Diagnostic Logging (insert inside training loop)]
# === Diagnostic Logging (insert inside training loop) ===
if epoch % R == 0:                                  # Rebuild cycle hook
    nnz = adj_matrix._nnz() if adj_matrix.is_sparse else (adj_matrix != 0).sum().item()
    logger.info(f"[epoch {epoch}] NNZ={nnz} | avg_deg={nnz / N:.2f}")

with torch.no_grad():                               # Tanh saturation probe
    sat = (torch.tanh(model.AVRF_image).abs() > 0.95).float().mean().item()
    sat_t = (torch.tanh(model.AVRF_text).abs() > 0.95).float().mean().item()
    logger.info(f"[epoch {epoch}] tanh_sat: img={sat:.3f} txt={sat_t:.3f}")
\end{lstlisting}

\begin{table}[htbp]
\centering
\renewcommand{\arraystretch}{1.3}
\begin{tabularx}{\textwidth}{|l|X|}
\hline
\textbf{Code / Command} & \textbf{Purpose} \\
\hline
\texttt{if epoch \% R == 0} & Only logs \texttt{NNZ} at rebuild cycles (avoiding overhead in intermediate epochs) aligning with Pattern B'. \\
\hline
\texttt{adj\_matrix.\_nnz()} & Standard PyTorch sparse tensor API; fallback to \texttt{(...!= 0).sum()} if using dense format to ensure safety. \\
\hline
\texttt{avg\_deg = nnz / N} & Logs the average degree per node—helping reviewers quickly verify $\approx K$ (proving \texttt{NNZ} is anchored). \\
\hline
\texttt{with torch.no\_grad()} & Avoids building unnecessary computational gradient graphs during probes. \\
\hline
\texttt{torch.tanh(...).abs() > 0.95} & Precisely defines the "tanh saturation rate"—the percentage of logits with $|x| > 0.95$. \\
\hline
\end{tabularx}
\end{table}

\subsubsection{Slim Momentum Design}
To minimize VRAM consumption, the Momentum Encoder applies strict Exponential Moving Average (EMA) smoothing only to the calibrated representations $h_{cal} \in \mathbb{R}^{d=64}$, rather than the entire raw multi-modal feature table. This ``Slim Momentum'' design helps reduce the auxiliary memory footprint by approximately 98\%.

\subsection{Engineering Bug Standardization and cuFFT Cache Control}
The definitive solution for the hidden cache memory leak is completely disabling the cuFFT plan cache via a static attribute, instead of abusing \texttt{empty\_cache}:
\begin{lstlisting}[language=Python]
# Static configuration immediately upon model initialization
torch.backends.cuda.cufft_plan_cache[device.index].max_size = 1
\end{lstlisting}

\section{Experimental Strategy, Reliability Evaluation, and Robustness}

\noindent \textbf{Phase 1 — Quick Win ($<$ 1 day) Execution Plan:} \\
Based on the latest review panel analysis, this blueprint directs all resources solely to Phase 1 to achieve an immediate recovery effect (Phases 2 and 3 are reserved for future deployment). Phase 1 consists of 3 highest-priority action items:
\begin{enumerate}
    \item \textbf{Audit eval protocol:} This is the highest priority with zero cost. A comprehensive cross-review of the source code is required to ensure the evaluation protocol perfectly matches the original MMHCL baseline. Any discrepancy here could potentially explain the entire 10.7\% Recall performance gap.
    \item \textbf{Static $\tau$ sweep:} Fix $\tau = 0.3$ to thoroughly eliminate the $\tau$ saturation bug at 0.0909, thereby restoring the diversity of the embedding space and improving coverage.
    \item \textbf{AVRF ablation experiment:} Temporarily disable the AVRF filter to preserve useful sparse signals from being over-attenuated in the multi-modal space.
\end{enumerate}

\noindent \textbf{Modern Baselines:} \\
To ensure fairness against State-of-the-Art (SOTA) models introduced after MMHCL, the evaluation space incorporates modern baselines: FREEDOM (SIGIR'23), MMSSL, LGMRec, and HPMRec.

\noindent \textbf{Bayesian Hyperparameter Optimization (Bayesian HPO):} \\
The large-scale optimization process will migrate to Bayesian HPO (Optuna/BOHB). (\textit{Note: During Phase 1, the static $\tau$ sweep will temporarily substitute this optimization process for the temperature parameter to deliver immediate results}).

\section{System Efficiency Report and Performance Expectations}

\noindent \textbf{Phase 1 Expected Gain:} \\
Asymmetric analyses indicate that the current DAMPS-MMHCL configuration suffers only a 2.9\% decline in NDCG@20 but a staggering 10.7\% drop in Recall@20. This asymmetric phenotype clearly demonstrates that the system is losing \textbf{coverage} of candidates (long-tail items) rather than failing at ranking.

For Phase 1 deployment results: Assuming a perfect match in the evaluation protocol, the combined remediation of a static $\tau$ and disabling AVRF ($\tau$+AVRF fix) is expected to yield a combined recovery and growth of \textbf{+5\% to +8\% Recall@20} on Amazon Clothing, nearly closing the current gap against the original baseline. The projected NDCG gains (+1\% to +4\%) across all datasets remain highly credible.

\subsection{System Efficiency Pilot Protocol}
All experiments utilize \texttt{bfloat16} mixed precision (natively supported on the RTX 5090 Ada Lovelace architecture), reducing actual wall-clock runtime by 30-40\% and VRAM usage by 40\%.

\begin{table}[htbp]
\centering
\renewcommand{\arraystretch}{1.2}
\begin{tabular}{|l|c|c|c|c|}
\hline
\textbf{Configuration} & \textbf{Rebuild R} & \textbf{Peak VRAM (GB)} & \textbf{Time/epoch (s)} & \textbf{Total Time (250e)} \\
\hline
MMHCL Baseline (w/o DAMPS) & - & 8.45 $\pm$ 0.12 & 12.04 $\pm$ 0.3 & 0.83 h \\
\hline
DAMPS-MMHCL, dense $N \times N$ matrix & 1 & 28.51 $\pm$ 0.45 & 25.10 $\pm$ 0.8 & 1.74 h \\
\hline
DAMPS-MMHCL, chunked & 1 & 11.23 $\pm$ 0.15 & 16.52 $\pm$ 0.4 & 1.15 h \\
\hline
DAMPS-MMHCL, chunked + FAISS & 1 & 11.48 $\pm$ 0.18 & 15.75 $\pm$ 0.3 & 1.10 h \\
\hline
DAMPS-MMHCL, chunked & \textbf{5} & \textbf{10.75 $\pm$ 0.11} & \textbf{12.96 $\pm$ 0.2} & \textbf{0.90 h} \\
\hline
DAMPS-MMHCL, chunked & 10 & 10.58 $\pm$ 0.10 & 12.45 $\pm$ 0.2 & 0.87 h \\
\hline
\end{tabular}
\end{table}

\section{Ablation Study Plan}
The ablation plan has been upgraded to encompass the anti-fragility properties of the system and to validate the Phase 1 strategy:
\begin{itemize}
    \item \textbf{MMHCL Baseline:} The fundamental data anchor point.
    \item \textbf{Graph Rebuild Frequency Ablation (Ablation over R):} Analyzes the correlation between time cost and performance scores.
    \item \textbf{Momentum Encoder Ablation (EMA on/off):} Examines the risk of signal attenuation and proves the stabilizing role of the feature smoothing structure.
    \item \textbf{Phase 1 - AVRF Ablation:} Completely disables the AVRF filter to measure the recovered coverage on the Amazon Clothing dataset.
    \item \textbf{+ Metadata-Aware APC Only:} Isolates the effect of the von Mises MLE estimation combined with static metadata categories.
    \item \textbf{+ Residual IMCF Only:} Analyzes the reduction of ``phantom edges'' based on the ASC coherence coefficient computed in a residual format.
    \item \textbf{Full DAMPS-MMHCL:} The convergent proposed model featuring Soft Residual-Routing.
    \item \textbf{Spectral Empirical Evaluation on Unordered Data (Permutation-FFT Test):} A direct comparison for the transformations involving: (a) 1D FFT, (b) DCT-II, (c) learnable orthogonal transformation, (d) spatial baseline, and (e) randomly permuted FFT. If the statistical gap between (a) and (e) is $< 1\%$, the system will switch to DCT-II.
\end{itemize}

\section{Conclusion}
The design blueprint for upgrading the MMHCL architecture via DAMPS has been strategically condensed to strike directly at \textbf{Phase 1 — Quick Win}.

By actively acknowledging and directly analyzing the asymmetry of the Recall/NDCG metrics, the current framework strictly prioritizes coverage recovery by: (1) Comprehensively auditing the evaluation protocol, (2) deploying a static InfoNCE temperature sweep fixed at $\tau = 0.3$ to break the embedding collapse caused by saturated gradients, and (3) temporarily disabling the AVRF filter to preserve invaluable sparse signals. Evaluations based on accurate baseline constraints and a practical Pilot benchmark table clearly delineate the system's smooth operational capabilities. With this focused pivot, the work is fully primed to deliver an immediate +5\%\textasciitilde8\% Recall@20 gain, firmly paving the way before advancing into Phase 2 and 3 optimizations.

\end{document}
